\addcontentsline{toc}{section}{\texorpdfstring{\MakeUppercase{List of Symbols}}{List of Symbols}}

\begin{center}
LIST OF SYMBOLS
\end{center}

\renewcommand{\arraystretch}{1.15}
% 35 entries per table!

\begin{center}
\begin{tabular}{>{$}l<{$} p{3.7in}}
\toprule
\text{Symbol} & \text{Meaning} \\
\midrule

q = M_2/M_1,\; M_1 > M_2 & mass ratio \\
M_1                      & mass of more massive binary component \\
M_2                      & mass of less massive binary component \\
\Xeff                    & effective spin \\
\tauAGN                  & AGN lifetime \\
\Rg = \frac{GM}{c^2}     & gravitational radius \\
G                        & gravitational constant \\
M                        & mass \\
c                        & speed of light \\
\vec{\chi}_{1,2}         & dimensionless spin vectors \\
\Lbin                    & binary orbital angular momentum \\
\Ldisk                   & disk angular momentum \\
\Lbh                     & black hole orbital angular momentum \\
\gamma                   & index of the mass power-law distribution \\
\Msun                    & solar mass \\
\Mbh                     & supermassive black hole mass \\
\mu_\chi                 & mean spin \\
\sigma_\chi              & spin standard deviation \\
\Rout                    & disk outer radius \\
H                        & disk scale height \\
e                        & orbital eccentricity \\
e_{\rm 0,max}            & maximum initial eccentricity \\
e_{\rm crit}             & eccentricity threshold where $a$ can evolve \\
a                        & orbital semi-major axis \\
\Rhill = a_1 (M_1/3\Mbh)^{1/3} & Hill radius \\
M_{\rm NSC}              & nuclear star cluster mass \\
N_{\rm BH}               & number of black holes \\
N_*                      & number of stars \\
\Rcrit                   & break radius in the NSC number density profile \\
\eta                     & power-law index for NSC profiles \\
n                        & number density \\
\fret                    & fraction of retrograde binaries \\

\end{tabular}
\end{center}


\begin{center}
\begin{tabular}{>{$}l<{$} p{3.7in}}
\toprule
\text{Symbol} & \text{Meaning~~~~~~~~~~~~~~~~~~~~~~~~~~~~~~~~~~~~~~~~~~~~~\textit{continued}} \\
\midrule

\mathcal{R}              & BBH merger rate \\
\varepsilon_{\rm SN}     & SN energy density \\
P_D(r)                   & disk gas pressure at radius r \\
T                        & gas temperature \\
k                        & Boltzmann's constant \\
\C                       & SN breakout parameter \\
\rho                     & gas volume density \\
\rho_0                   & disk midplane density \\
\Esn                     & SN total kinetic energy \\
\Rsn                     & SN shock radius \\
\mProt                   & proton mass \\
\Omega_0                 & Keplerian orbital frequency \\
q_\text{shear}           & logarithmic shearing parameter \\
E                        & total energy density \\
\cs                      & sound speed \\
\rhoej                   & SN ejecta density \\
\Mej                     & SN ejecta mass \\
\tpds                    & PDS transition time \\
\Rpds                    & PDS transition radius \\
\fkinetic                & kinetic energy fraction \\
\Rm                      & muffling radius \\
\Rs                      & Schwarzschild radius \\
T(z)                     & vertical temperature profile \\
\Tmid                    & disk midplane temperature \\
\Tiso                    & isothermal temperature cap \\
\Tfloor                  & minimum temperature floor \\
\Teff                    & disk effective temperature \\
H_T                      & temperature scale height \\
\tau                     & optical depth \\
\theta = 5040/T          & thermal parameter \\
\kappa_\lambda           & wavelength-dependent absorption opacity \\
\sigma_\lambda           & wavelength-dependent scattering opacity \\
I_\lambda                & specific intensity \\


\end{tabular}
\end{center}


\begin{center}
\begin{tabular}{>{$}l<{$} p{3.7in}}
\toprule
\text{Symbol} & \text{Meaning~~~~~~~~~~~~~~~~~~~~~~~~~~~~~~~~~~~~~~~~~~~~~\textit{continued}} \\
\midrule

S_\lambda                & source function \\
B_\lambda                & Planck function \\
J_\lambda                & mean intensity \\
\omega_\lambda           & single-scattering albedo \\
\Iplus                   & upward intensity \\
\Iminus                  & downward intensity \\
U = (I^+ + I^-)/2        & mean intensity (Feautrier) \\
V = (I^+ - I^-)/2        & flux (Feautrier) \\
\alpha_\lambda           & total extinction coefficient \\
F_\lambda                & flux density \\
\Hminus                  & negative hydrogen ion \\
\Hplus                   & ionized hydrogen \\
\nElec                   & electron number density \\
\nI                      & ion number density \\
\gff                     & free-free Gaunt factor \\
g_{\rm bf}               & bound-free Gaunt factor \\
P_{\rm e}                & partial electron pressure \\
\chi_n                   & excitation potential for level n \\
I = 13.598\ {\rm eV}     & hydrogen ionization energy \\
\sigma_0                 & absorption cross-section normalization \\
\lambda                  & wavelength \\
\nu                      & frequency \\
h                        & Planck constant (distinct from disk scale height H) \\
\me                      & electron mass \\
\chi_{\rm p}             & in-plane spin component \\
\mathcal{M}_{\rm ch}     & binary chirp mass \\
\mu_{\rm BBH}            & BBH reduced mass \\
D                        & distance to source \\
\nu_{\rm GW}             & gravitational wave frequency \\
h_{\rm char}             & characteristic GW strain \\
h_\text{str}             & GW strain \\
a_{\rm BBH}              & BBH semi-major axis \\
M_{\rm BBH}              & BBH total mass \\

\end{tabular}
\end{center}


\begin{center}
\begin{tabular}{>{$}l<{$} p{3.7in}}
\toprule
\text{Symbol} & \text{Meaning~~~~~~~~~~~~~~~~~~~~~~~~~~~~~~~~~~~~~~~~~~~~~\textit{continued}} \\
\midrule

\dot{\nu}_{\rm GW}       & GW frequency time derivative \\
N                        & number of GW cycles \\
t_{\rm damp}             & orbital eccentricity damping timescale \\
h = H/r                  & disk aspect ratio (distinct from Planck constant h) \\
n_{\rm AGN}              & AGN number density \\
\alpha                   & disk viscosity parameter \\
p(R)                     & probability density over radius \\
\Delta R                 & thickness of a spherical shell \\

\end{tabular}
\end{center}
